\documentclass[10pt]{beamer}
\usetheme{metropolis}
\usepackage{graphicx}
\usepackage{listings}
\usepackage{booktabs}
\usepackage{xcolor}
\definecolor{codebg}{HTML}{F5F5F5}
\definecolor{codegreen}{HTML}{2E7D32}
\definecolor{codegray}{HTML}{757575}
\definecolor{codeblue}{HTML}{1565C0}
\lstdefinestyle{rstyle}{language=R,backgroundcolor=\color{codebg},basicstyle=\ttfamily\scriptsize,
  keywordstyle=\color{codeblue}\bfseries,stringstyle=\color{codegreen},
  commentstyle=\color{codegray}\itshape,breaklines=true,frame=single,
  rulecolor=\color{codegray!30},showstringspaces=false,tabsize=2,
  morekeywords={library,filter,select,mutate,group_by,summarise,pivot_longer,
    left_join,ggplot,aes,geom_col,coord_flip,theme_minimal,labs}}
\lstdefinestyle{pythonstyle}{language=Python,backgroundcolor=\color{codebg},basicstyle=\ttfamily\scriptsize,
  keywordstyle=\color{codeblue}\bfseries,stringstyle=\color{codegreen},
  commentstyle=\color{codegray}\itshape,breaklines=true,frame=single,
  rulecolor=\color{codegray!30},showstringspaces=false,tabsize=4,
  morekeywords={import,from,as,pd,np,plt,sns,query,assign,sort_values,groupby,
    agg,value_counts,nlargest,dropna,fillna,to_numeric,to_datetime,
    melt,pivot_table,merge,drop_duplicates,str,replace,rename,astype,
    explode,where,between,pipe,lambda,apply,map,nsmallest,DataFrame}}
\title{Data Wrangling for Visualization (Python)}
\subtitle{Workshop 4 --- Module 3}
\author{Prof.Asoc.\ Endri Raco}
\institute{Data Visualization for Data Scientists \\ UNYT Tirana}
\date{2025/26}
\begin{document}

\begin{frame}
  \titlepage
\end{frame}

\begin{frame}{Module Roadmap}
  \begin{enumerate}
    \item The pandas pipeline: method chaining as Python's pipe
    \item R $\leftrightarrow$ Python Rosetta Stone: every dplyr verb mapped
    \item .assign(), .query(), .groupby().agg(): the core trio
    \item Reshaping with pd.melt() and .pivot\_table()
    \item Common cleaning patterns in pandas
    \item Full chain: one expression from CSV to chart
  \end{enumerate}
  \begin{exampleblock}{Learning Outcome}
    You will translate any R/dplyr wrangling pipeline into idiomatic
    pandas method chains, and produce the same ggplot-ready data
    structures using Python's .groupby().agg() and pd.melt().
  \end{exampleblock}
\end{frame}

\section{The pandas Pipeline}

\begin{frame}{Method Chaining: Python's Pipe}
  \begin{center}
    \includegraphics[width=0.95\textwidth]{figures_w04m03/pandas_pipeline}
  \end{center}
\end{frame}

\begin{frame}{Chaining vs Temporary Variables}
  \begin{center}
    \includegraphics[width=0.92\textwidth]{figures_w04m03/method_chaining}
  \end{center}
  \begin{alertblock}{Pro-Tip}
    Wrap the entire chain in parentheses so Python allows multi-line
    method calls without backslash continuation. The \texttt{.assign()}
    method creates columns inline (like dplyr's \texttt{mutate()}).
    Use \texttt{lambda d:} inside \texttt{.assign()} to reference
    the DataFrame being transformed.
  \end{alertblock}
\end{frame}

\section{R $\leftrightarrow$ Python Rosetta Stone}

\begin{frame}{Every dplyr Verb $\rightarrow$ pandas Equivalent}
  \begin{center}
    \includegraphics[width=0.95\textwidth]{figures_w04m03/rosetta}
  \end{center}
\end{frame}

\begin{frame}[fragile]{.query() + Indexing: Subset Rows and Columns}
\begin{lstlisting}[style=pythonstyle]
# .query(): string-based filtering (like dplyr::filter)
df.query("Rating > 4.0 and Type == 'Free'")
df.query("Category in ['GAME', 'FAMILY', 'TOOLS']")

# Boolean indexing: equivalent but more verbose
df[df["Rating"] > 4.0]
df[(df["Rating"] > 4) & (df["Type"] == "Free")]  # & not 'and'

# Column selection (like dplyr::select)
df[["App", "Category", "Rating", "Type"]]     # keep columns
df.drop(columns=["Reviews", "Installs"])       # drop columns
df.filter(like="Rating")                       # pattern match
df.select_dtypes(include="number")             # by type

# .query() with variables (use @)
min_rating = 4.0
df.query("Rating > @min_rating")               # reference variable
\end{lstlisting}
\end{frame}

\begin{frame}[fragile]{.assign(): Create and Transform Columns}
\begin{lstlisting}[style=pythonstyle]
# .assign(): chainable column creation (like dplyr::mutate)
(df
    .assign(
        # Type conversion
        Reviews=lambda d: pd.to_numeric(d["Reviews"], errors="coerce"),
        # Log transform
        log_reviews=lambda d: np.log10(d["Reviews"] + 1),
        # String cleaning
        Size_MB=lambda d: d["Size"].str.replace("M", "").pipe(
            pd.to_numeric, errors="coerce"),
        # Date parsing
        date=lambda d: pd.to_datetime(d["Last Updated"],
            format="%B %d, %Y", errors="coerce"),
        year=lambda d: d["date"].dt.year,
        # Conditional recode (like case_when)
        price_cat=lambda d: np.select(
            [d["Price"] == 0, d["Price"] < 5, d["Price"] < 20],
            ["Free", "Under $5", "$5-$20"], default="$20+")
    )
)
\end{lstlisting}
\end{frame}

\begin{frame}[fragile]{.groupby().agg(): Aggregate per Group}
\begin{lstlisting}[style=pythonstyle]
# Named aggregation (cleanest syntax, pandas 0.25+)
(df
    .groupby("Category")
    .agg(
        n=("Rating", "count"),
        mean_r=("Rating", "mean"),
        median_r=("Rating", "median"),
        std_r=("Rating", "std"),
        q25=("Rating", lambda x: x.quantile(0.25))
    )
    .sort_values("mean_r", ascending=False)
)

# Quick counting (like dplyr::count)
df["Category"].value_counts()             # Series, sorted desc
df.groupby("Category").size()             # same but groupby-based

# Top-n (like dplyr::slice_max)
df.nlargest(10, "Rating")                 # top 10 by Rating
summary.nlargest(8, "n")                  # top 8 by count
\end{lstlisting}
\end{frame}

\section{Reshaping}

\begin{frame}{pd.melt() $\leftrightarrow$ .pivot\_table()}
  \begin{center}
    \includegraphics[width=0.92\textwidth]{figures_w04m03/melt_pivot}
  \end{center}
  \begin{exampleblock}{Data Insight}
    seaborn and plotnine require long (tidy) format, just like ggplot2.
    \texttt{pd.melt()} converts wide$\rightarrow$long.
    \texttt{.pivot\_table()} converts long$\rightarrow$wide (for summary
    tables). The \texttt{.explode()} method handles multi-valued cells
    (e.g., Netflix's genre column).
  \end{exampleblock}
\end{frame}

\begin{frame}[fragile]{Reshaping in Practice}
\begin{lstlisting}[style=pythonstyle]
# WIDE -> LONG (for seaborn / plotnine)
df_long = pd.melt(
    df_wide,
    id_vars=["region"],            # keep these columns
    var_name="year",                # new col for former col names
    value_name="revenue"            # new col for values
)
df_long["year"] = df_long["year"].astype(int)  # string -> int

# LONG -> WIDE (for summary tables)
df_wide = df_long.pivot_table(
    index="region",
    columns="year",
    values="revenue",
    aggfunc="sum"                   # handles duplicates
)

# EXPLODE: one row per value in a list column
netflix["genres"] = netflix["listed_in"].str.split(", ")
netflix_exploded = netflix.explode("genres")  # like separate_rows()
\end{lstlisting}
\end{frame}

\section{Cleaning Patterns}

\begin{frame}{Six Common Cleaning Patterns}
  \begin{center}
    \includegraphics[width=0.92\textwidth]{figures_w04m03/cleaning_patterns}
  \end{center}
\end{frame}

\begin{frame}[fragile]{Cleaning the Google Play Store Dataset}
\begin{lstlisting}[style=pythonstyle]
apps_clean = (
    pd.read_csv("googleplaystore.csv")
    .drop_duplicates(subset="App", keep="first")
    .assign(
        Reviews=lambda d: pd.to_numeric(d["Reviews"], errors="coerce"),
        Installs=lambda d: (d["Installs"]
            .str.replace(r"[+,]", "", regex=True)
            .pipe(pd.to_numeric, errors="coerce")),
        Size_MB=lambda d: (d["Size"]
            .str.replace("M", "").str.replace("k", "e-3")
            .pipe(pd.to_numeric, errors="coerce")),
        Price=lambda d: (d["Price"]
            .str.replace("$", "", regex=False)
            .pipe(pd.to_numeric, errors="coerce")),
        date=lambda d: pd.to_datetime(
            d["Last Updated"], format="%B %d, %Y", errors="coerce")
    )
    .dropna(subset=["Rating"])
    .query("Type in ['Free', 'Paid']")
    .query("Rating >= 1 and Rating <= 5")
)
\end{lstlisting}
\end{frame}

\section{Joins}

\begin{frame}[fragile]{pd.merge(): Combining DataFrames}
\begin{lstlisting}[style=pythonstyle]
# LEFT JOIN (most common for viz)
result = pd.merge(
    apps, lookup_table,
    on="Category",           # join key
    how="left"               # keep all apps rows
)

# INNER JOIN
pd.merge(apps, reviews, on="App", how="inner")

# ANTI JOIN (rows in apps with NO match in reviews)
anti = apps[~apps["App"].isin(reviews["App"])]
# Or with merge:
merged = pd.merge(apps, reviews, on="App", how="left", indicator=True)
anti = merged.query("_merge == 'left_only'").drop(columns="_merge")

# JOIN TYPES: how = "left", "right", "inner", "outer"
\end{lstlisting}
  \begin{alertblock}{Pro-Tip}
    pandas uses \texttt{pd.merge()} or \texttt{df.merge()} with
    \texttt{how=} parameter. For anti-joins, use \texttt{indicator=True}
    then filter on \texttt{\_merge == 'left\_only'}, or simply use
    \texttt{\textasciitilde{}df["key"].isin(other["key"])}.
  \end{alertblock}
\end{frame}

\section{Full Chain}

\begin{frame}{One Expression: CSV $\rightarrow$ Chart}
  \begin{center}
    \includegraphics[width=0.82\textwidth]{figures_w04m03/full_chain}
  \end{center}
\end{frame}

\section{Complete Examples}

\begin{frame}[fragile]{R: Pipe-to-Plot (Reference)}
  \begin{columns}[T]
    \begin{column}{0.52\textwidth}
\begin{lstlisting}[style=rstyle]
# R equivalent for comparison
apps_clean |>
  group_by(Category) |>
  summarise(
    mean_r = mean(Rating),
    n = n(),
    .groups = "drop") |>
  slice_max(n, n = 8) |>
  mutate(Category =
    fct_reorder(Category, mean_r),
    hl = mean_r == max(mean_r)) |>
  ggplot(aes(x = Category,
    y = mean_r, fill = hl)) +
  geom_col(width = 0.5,
    show.legend = FALSE) +
  scale_fill_manual(values =
    c("TRUE" = "#1565C0",
      "FALSE" = "#BBBBBB")) +
  coord_flip() +
  theme_minimal() +
  labs(title = "R: pipe-to-plot")
\end{lstlisting}
    \end{column}
    \begin{column}{0.45\textwidth}
      \includegraphics[width=\textwidth]{figures_w04m03/r_output}
    \end{column}
  \end{columns}
\end{frame}

\begin{frame}[fragile]{Python: Chain-to-Plot}
  \begin{columns}[T]
    \begin{column}{0.52\textwidth}
\begin{lstlisting}[style=pythonstyle]
# Python chain equivalent
summary = (apps_clean
    .groupby("Category")
    .agg(
        mean_r=("Rating", "mean"),
        n=("Rating", "count"))
    .nlargest(8, "n")
    .sort_values("mean_r"))

fig, ax = plt.subplots(
    figsize=(6, 5))
colors = ["#BBBBBB"] * len(summary)
colors[-1] = "#2E7D32"  # accent

ax.barh(summary.index,
    summary["mean_r"],
    color=colors, height=0.5)

for i, (idx, row) in enumerate(
    summary.iterrows()):
    ax.text(row["mean_r"] + 0.005,
        i, f"{row['mean_r']:.2f}",
        va="center", fontsize=7)

ax.set_xlim(3.8, 4.4)
ax.set_title("Python: chain-to-plot")
\end{lstlisting}
    \end{column}
    \begin{column}{0.45\textwidth}
      \includegraphics[width=\textwidth]{figures_w04m03/py_output}
    \end{column}
  \end{columns}
\end{frame}

\section{Synthesis}

\begin{frame}{Key Takeaways}
  \begin{enumerate}
    \item \textbf{Method chaining} (parenthesised multi-line) is pandas'
          equivalent of R's \texttt{|>} pipe. Wrap in \texttt{(...)}.
    \item \textbf{.assign()} = \texttt{mutate()};
          \textbf{.query()} = \texttt{filter()};
          \textbf{.groupby().agg()} = \texttt{group\_by() + summarise()}.
    \item \textbf{pd.melt()} = \texttt{pivot\_longer()};
          \textbf{.pivot\_table()} = \texttt{pivot\_wider()};
          \textbf{.explode()} = \texttt{separate\_rows()}.
    \item Common cleaning: \texttt{pd.to\_numeric(errors='coerce')},
          \texttt{.str.replace()}, \texttt{pd.to\_datetime()},
          \texttt{np.where()} / \texttt{np.select()}.
    \item \textbf{Full chain}: one parenthesised expression from
          \texttt{pd.read\_csv()} to \texttt{.plot.barh()} --- zero
          intermediate variables.
    \item For anti-joins: \texttt{$\sim$df["key"].isin(other["key"])}
          or \texttt{merge(indicator=True)}.
  \end{enumerate}
\end{frame}

\begin{frame}{Essential Readings}
  \begin{itemize}
    \item McKinney, W. (2022). \emph{Python for Data Analysis}, 3rd ed.,
          Chapters 5, 7--8 (pandas, Data Wrangling).
    \item pandas documentation: Method chaining style guide.
          \texttt{https://pandas.pydata.org/docs/}
    \item VanderPlas, J. (2016). \emph{Python Data Science Handbook},
          Chapter 3 (pandas).
  \end{itemize}
  \vspace{6pt}
  \textbf{Next Module}: Missing Data Visualization (W04-M04)
\end{frame}

\end{document}
